\section{Translated dyadic neighborhoods}
\label{sec:area_law_dyadic_neighborhoods}

\begin{definition}[Half-open dyadic cells]
    \label{def:area_law_dyadic_cells}
    \lean{TNLean.PEPS.AreaLaw.Geometry.dyadicCell,
        TNLean.PEPS.AreaLaw.Geometry.dyadicCellIndex,
        TNLean.PEPS.AreaLaw.Geometry.dyadicParent}
    \leanok
    Fix an origin $o\in\mathbb R^2$. For every integer $k\ge0$ and index
    $z\in\mathbb Z^2$, define the half-open cell and its index map by
    \begin{align}
        C_{k,z} & =o+2^k\bigl(z+[0,1)^2\bigr),                         \\
        c_k(x)  & =\left(\left\lfloor\frac{x_1-o_1}{2^k}\right\rfloor,
        \left\lfloor\frac{x_2-o_2}{2^k}\right\rfloor\right),           \\
        p(z)    & =\left(\left\lfloor\frac{z_1}{2}\right\rfloor,
        \left\lfloor\frac{z_2}{2}\right\rfloor\right).
    \end{align}
    The map $p$ gives the index of a cell's parent. The arbitrary origin
    includes the translated origin used in \cite[Section~11]{OpenAI2026AreaLaw}.
\end{definition}

\begin{theorem}[Cell indices and parents]
    \label{thm:area_law_dyadic_indices}
    \lean{TNLean.PEPS.AreaLaw.Geometry.mem_dyadicCell_iff,
        TNLean.PEPS.AreaLaw.Geometry.dyadicCellIndex_succ,
        TNLean.PEPS.AreaLaw.Geometry.dyadicParent_abs_sub_le,
        TNLean.PEPS.AreaLaw.Geometry.dyadicParent_mem_ambientDilation_image,
        TNLean.PEPS.AreaLaw.Geometry.mem_dyadicChildren_iff}
    \leanok
    \uses{def:area_law_dyadic_cells,def:area_law_cell_counts,def:area_law_template}
    Every point lies in the unique cell indexed by $c_k(x)$, and passing to
    parents agrees with increasing the scale:
    \begin{align}
        x\in C_{k,z} & \quad\Longleftrightarrow\quad c_k(x)=z, \\
        c_{k+1}(x)   & =p(c_k(x)).
    \end{align}
    For every integer $C_0\ge0$, two integer indices at sup distance at most
    $C_0$ have parents at sup distance at most $C_0$. Consequently $p$ maps
    the integer $C_0$-neighborhood of a finite index set into the
    $C_0$-neighborhood of its parent image. An index is a child
    of a member of $P$ precisely when its parent belongs to $P$.
\end{theorem}
\begin{proof}\leanok
    The defining inequalities of the floor function give the cell membership
    criterion. The identity $\lfloor t/2\rfloor=\lfloor\lfloor t\rfloor/2\rfloor$
    gives the parent formula. Integer division by two preserves the stated
    coordinate bounds, including negative indices. Finally the remainder modulo
    two is either zero or one, identifying the unique child in each coordinate.
\end{proof}

\begin{definition}[Occupied indices, neighborhoods and layers]
    \label{def:area_law_dyadic_layers}
    \lean{TNLean.PEPS.AreaLaw.Geometry.occupiedCellIndices,
        TNLean.PEPS.AreaLaw.Geometry.dyadicNeighborhood,
        TNLean.PEPS.AreaLaw.Geometry.dyadicLayerIndices,
        TNLean.PEPS.AreaLaw.Geometry.dyadicLayer}
    \leanok
    \uses{def:area_law_dyadic_cells,def:area_law_cell_counts,def:area_law_template}
    Let $Z\subset\mathbb Z^2$ be finite and let $C_0\ge0$ be an integer.
    For each integer $k\ge0$, set
    \begin{align}
        \mathcal Z_k & =\{c_k(z):z\in Z\},                          \\
        P_k          & =(\mathcal Z_k+[-C_0,C_0]^2)\cap\mathbb Z^2, \\
        N_k          & =\bigcup_{z\in P_k}C_{k,z},                  \\
        D_k          & =N_{k+1}\setminus N_k,                       \\
        I_k          & =\operatorname{Ch}(P_{k+1})\setminus P_k.
    \end{align}
    These are actual subsets of the plane and finite sets of indices; no
    nesting or distance bound is assumed in their definitions.
    The picture below shows one endpoint with $C_0=1$: the layer $D_k$ is
    the ring of scale-$k$ cells that lie in $N_{k+1}$ but not in $N_k$.
    % Source: 10-geometry.tex, lines 148--200 (N_k, D_k = N_{k+1} minus N_k);
    % no source figure. Specialization drawn: Z is one endpoint and C_0 = 1.
    % Ink-to-index: every dot is one half-open scale-k cell C_{k,z}; the red dot
    % is the occupied index c_k(z) in Z_k. The pink enclosure is the 3-by-3
    % block P_k, so its union is N_k. Each dashed square is one parent cell
    % C_{k+1,p} of P_{k+1}, i.e. its four scale-k children; the nine of them
    % make N_{k+1}, marked by the bracket. The tinted ring is the layer D_k,
    % the union of the cells in I_k = Ch(P_{k+1}) minus P_k. N_k is not aligned
    % with the parent squares, which is why D_k cuts through four of them.
    % No strings, legs or contractions: this is a geometry picture with empty
    % boundary signature. Not to scale; the plane continues past the frame.
    \begin{center}
        \begin{tenkz}[size=s, rows={ket,ket,ket,ket,ket,ket,ket,ket}, cols=8, bonds=none]
            \tndeclare{species}{endpoint}{hue=source:red}
            \tn[at=(5,5), skin=dot, species=endpoint]{}
            \tnmark[form=enclosure, species=secondary, tint]{(2,2) .. (7,7) - (4,4) .. (6,6)}{}
            \tnmark[form=enclosure, species=selected, tint]{(4,4) .. (6,6)}{}
            \tnmark[form=enclosure, species=complement]{(2,2) .. (3,3)}{}
            \tnmark[form=enclosure, species=complement]{(2,4) .. (3,5)}{}
            \tnmark[form=enclosure, species=complement]{(2,6) .. (3,7)}{}
            \tnmark[form=enclosure, species=complement]{(4,2) .. (5,3)}{}
            \tnmark[form=enclosure, species=complement]{(4,4) .. (5,5)}{}
            \tnmark[form=enclosure, species=complement]{(4,6) .. (5,7)}{}
            \tnmark[form=enclosure, species=complement]{(6,2) .. (7,3)}{}
            \tnmark[form=enclosure, species=complement]{(6,4) .. (7,5)}{}
            \tnmark[form=enclosure, species=complement]{(6,6) .. (7,7)}{}
            \tnmark[form=label, label pos=e]{midway (2,2) and (3,2)}{$D_k$}
            \tnmark[form=label, label pos=e]{midway (4,4) and (5,4)}{$N_k$}
            \tnmark[form=bracket, label pos=e]{(2,2) .. (7,7)}{$N_{k+1}$}
        \end{tenkz}
    \end{center}
\end{definition}

\begin{theorem}[Nested neighborhoods and layer cells]
    \label{thm:area_law_dyadic_nesting}
    \lean{TNLean.PEPS.AreaLaw.Geometry.occupiedCellIndices_succ,
        TNLean.PEPS.AreaLaw.Geometry.mem_dyadicNeighborhood_iff,
        TNLean.PEPS.AreaLaw.Geometry.dyadicNeighborhood_subset_succ,
        TNLean.PEPS.AreaLaw.Geometry.dyadicLayer_eq_iUnion}
    \leanok
    \uses{def:area_law_dyadic_layers,thm:area_law_dyadic_indices}
    For every origin, finite $Z$, integer $C_0\ge0$, and integer $k\ge0$,
    \begin{align}
        \mathcal Z_{k+1} & =p(\mathcal Z_k),                            \\
        x\in N_k         & \quad\Longleftrightarrow\quad c_k(x)\in P_k, \\
        N_k              & \subseteq N_{k+1},                           \\
        D_k              & =\bigcup_{z\in I_k}C_{k,z}.
    \end{align}
    The statements include the empty set of endpoints and radius zero.
\end{theorem}
\begin{proof}\leanok
    Apply the parent identity to every point of $Z$. The cell membership
    criterion identifies membership in $N_k$ with membership of its index in
    $P_k$. The parent preserves the integer radius, proving nesting.
    A point belongs to $N_{k+1}$ precisely when its scale-$k$ index is a child
    of an index in $P_{k+1}$. Removing $N_k$ therefore removes precisely the
    indices in $P_k$, which gives the last equality.
\end{proof}

\begin{theorem}[Uniform number of layer cells]
    \label{thm:area_law_dyadic_layer_count}
    \lean{TNLean.PEPS.AreaLaw.Geometry.card_dyadicLayerIndices_le,
        TNLean.PEPS.AreaLaw.Geometry.card_dyadicLayerIndices_boundary_le}
    \leanok
    \uses{thm:area_law_cell_counts,thm:area_law_dyadic_nesting}
    For every origin, finite $Z$, integer $C_0\ge0$, and integer $k\ge0$,
    \begin{align}
        |I_k| & \le4(2C_0+1)^2|Z|.
    \end{align}
    When $Z$ is the set of endpoints of the crossing edges of a cut in a
    finite induced domain, with $b=|\partial_\Lambda A|$, this gives
    \begin{align}
        |I_k| & \le8(2C_0+1)^2b.
    \end{align}
    The constant depends only on the chosen radius, independently of the
    origin, scale, domain and cut. These are the layer-cell estimates in
    \cite[Section~11]{OpenAI2026AreaLaw}.
\end{theorem}
\begin{proof}\leanok
    The occupied-index image has at most $|Z|$ elements. Its integer
    $C_0$-neighborhood has at most $(2C_0+1)^2|Z|$ elements, and each parent
    has exactly four children. Taking a set difference can only decrease the
    count. Each crossing edge contributes at most two endpoints, giving the
    second bound.
\end{proof}


\begin{theorem}[Exhaustion of the plane]
    \label{thm:area_law_dyadic_exhaustion}
    \lean{TNLean.PEPS.AreaLaw.Geometry.exists_mem_dyadicNeighborhood,
        TNLean.PEPS.AreaLaw.Geometry.iUnion_dyadicNeighborhood_eq_univ}
    \leanok
    \uses{def:area_law_dyadic_layers,thm:area_law_dyadic_indices}
    For every origin, nonempty finite endpoint set $Z$, and integer $C_0\ge2$,
    every point belongs to some $N_k$, with $k\ge0$. Consequently,
    \begin{align}
        \bigcup_{k\ge0}N_k & =\mathbb R^2.
    \end{align}
    This is the exhaustion assertion in \cite[Section~11]{OpenAI2026AreaLaw}.
\end{theorem}
\begin{proof}\leanok
    Choose $z\in Z$. For any $x\in\mathbb R^2$, take $2^k$ larger than the
    sum of the four absolute values of the coordinates of $x-o$ and $z-o$.
    Each translated coordinate divided by $2^k$ then lies in $(-1,1)$, so
    its floor is either $-1$ or $0$. The indices $c_k(x)$ and $c_k(z)$
    therefore differ by at most one in each coordinate. Since $c_k(z)$ is
    occupied and $C_0\ge2$, the index $c_k(x)$ belongs to $P_k$. Thus
    $x\in N_k$, as required.
\end{proof}


\begin{theorem}[Closures of cells and layers]
    \label{thm:area_law_dyadic_closures}
    \lean{TNLean.PEPS.AreaLaw.Geometry.closure_dyadicCell,
        TNLean.PEPS.AreaLaw.Geometry.closure_dyadicLayer_eq_iUnion}
    \leanok
    \uses{thm:area_law_dyadic_nesting}
    For every origin, nonnegative scale and cell index, the closure of a
    half-open dyadic cell is the corresponding closed rectangle. For every
    finite endpoint set and nonnegative integer radius, the actual layer satisfies
    \begin{align}
        \overline{C_{k,z}} & =o+2^k\bigl(z+[0,1]^2\bigr),           \\
        \overline{D_k}     & =\bigcup_{z\in I_k}\overline{C_{k,z}}.
    \end{align}
    These formulas give the closed regions used in the distance estimates of
    \cite[Section~11]{OpenAI2026AreaLaw}.
\end{theorem}
\begin{proof}\leanok
    The side length $2^k$ is positive, so the closure of each half-open coordinate
    interval is its closed interval. Closure commutes with finite products.
    The exact layer union is finite, and closure commutes with finite unions,
    which gives the second formula.
\end{proof}
